\section{The arithmetic meaning of a quantum prediction}\label{sec:arithmetic}
Consider a spin prepared in amplitudes $(z,w)$, sent through a Hadamard gate and then measured in the original basis. The first outcome has probability $|z+w|^2/2$. It matters whether a device adds the amplitudes before squaring. For $z=w=1/\sqrt2$ the first outcome is certain, although the incoherent mixture with the same initial diagonal has first-outcome probability $1/2$. The identity
\[
 |z+w|^2+|z-w|^2=2(|z|^2+|w|^2)
\]
certifies normalization without evaluating a square root. Clearing denominators turns this conservation law into an integer polynomial identity. The same principle of finite observation works for a wave packet with infinitely many momentum values, once a quantitative rule supplies its neglected norm. The significant mathematical work is to construct such rules and show that predictions depend only on the represented objects.

\subsection{Theory, codes and logic}
We use the ordinary iPRA of ProveIt \cite{proveit}. Its domain consists of natural numbers. The language has zero, successor, and a symbol for each primitive recursive (PR) definition. Its axioms give the defining equations, successor injectivity and separation from zero, and induction on quantifier-free formulas with free numerical parameters. The logic is intuitionistic first-order predicate logic with equality. Parameters, requested precisions, integer dimensions, time bounds and cutoff indices are all numerical. ``For every presented Hilbert vector'' is a metatheoretic schema: fix the PR symbols for its samplers, introduce any numerical parameters as free variables, and assert an implication from the stated first-order laws. Every explicit numerical family with one joint PR sampler gives a single parameterized instance. For an operation described as a supplied evaluator on effects or names, require that applying it to each joint PR family used in the proof gives a joint PR family, with the stipulated laws. An evaluator on arbitrary natural-number codes for PR programs would be a different and stronger assumption. A law supplied as a hypothesis remains a hypothesis until its instance has a derivation.

Here is one concrete choice of coding. Represent a signed integer by $(p,n)$ with value $p-n$, a rational by $(p,n,d)$ with value $(p-n)/(d+1)$, and a complex rational by two rationals. Equality and order reduce to natural-number comparisons by cross multiplication. A length and a nested pair of entries encode a finite list. For example $\pi(u,v)=(u+v)^2+u$ has projections found by two bounded loops up to the input code; give unused codes the default value zero. Reading a list, scanning a matrix, summing and multiplying by a length bound, concatenation and building a code as a loop accumulator are PR. A maximum, a least index within a stated bound, and a bounded universal test are PR (fold the conjunction of the individual Boolean answers). The same operations implement arbitrary rectangular shapes, including empty environments when an empty sum is appropriate. Positive denominators can always be retained by products, without reduction to lowest terms. Dimensions of physical state spaces are positive.

Write $\delta_k=2^{-k}$. A real or complex scalar is a fixed PR sampler of rational values $a_k$ satisfying $|a_k-a_l|\le\delta_k$ for $l\ge k$, with an integer magnitude bound. For complex values use the squared absolute value to check the inequality. Equality means $|a_k-b_k|\le2\delta_k$ for every $k$. The closed order on reals is $a_k\le b_k+2\delta_k$ for every $k$. An apartness or positivity witness gives a specified positive rational lower bound, possibly after computing more digits. Sample values that are certified to within $\delta_{k+2}$ of the intended answer at stage $k$ define a regular name. Section~\ref{sec:infinitehilbert} proves the corresponding results for vector names. An operator estimate $\|T\|\le b$ denotes an inequality on all its finite-vector tests. It never requires forming a real supremum.

Scalar addition and multiplication use the triangle inequality and the estimate
$|ab-a'b'|\le |a||b-b'|+|a-a'||b'|$. To compute a positive square root on $[0,B]$, approximate the input to $\delta_{2k+6}$, clamp the rational to $[0,B]$, and bisect the interval $[0,1+B]$ for its square root to width $\delta_{k+3}$. This costs at most $k+4+\lceil B\rceil$ steps. The bound $|\sqrt u-\sqrt v|^2\le|u-v|$ follows by factoring the difference of squares. This reasoning may first define square roots of rationals by the bisection itself, and then extend to names. Division by a real of magnitude at least $b>0$ is PR in its name and the rational witness, since input error $\eta\le b/2$ gives reciprocal error at most $2\eta/b^2$. For complex division use the conjugate divided by the squared modulus. These rules also yield absolute values and rational upper bounds for subsequent calculations. No sign decision on a general named real is used.

\subsection{From estimates to first-order arguments}
There are three recurring moves in the proofs.
\begin{enumerate}
\item A closed inequality at the limit is proved to arbitrary rational slack. If $r,s$ are rational observations and $r\le s+C2^{-m}$ for every $m$, a violation $r-s=u/v>0$ gives a contradictory finite test at $m=1+\lceil Cv\rceil$ for integers $u,v\ge1$. When $r,s$ themselves are named, first fix their observation stage and include its prescribed errors. Thus removal of slack uses only rational decidability and bounded integer arithmetic.
\item A limit is constructed only with a joint PR sampler and a PR Cauchy rate. Given the rate $m(k)$, take the $(k+3)$rd sample at index $m(k+3)$ (use running maxima for the rate). The three errors, from the first sample, the Cauchy comparison and the second sample, establish regularity. The same argument proves a uniform operator estimate by testing one finite input vector at a time, scaling the tolerance with its known bound.
\item For a loop of length $n$, fix its requested output tolerance before induction. Rounding each addition to error $\eta$ contributes at most $n\eta$. If a step amplifies error by $L\ge1$, the accumulated error with per-step residual $\eta$ is at most $nL^n\eta$ plus $L^n$ times the initial error. Prove the finite prefix inequality by quantifier-free induction on the actual rational arrays with these fixed parameters. For nonuniform loops replace the powers by products of the supplied integer step bounds. When a hypothesis is a law at every precision, instantiate it at the finitely many precisions needed by the computation. Bounded conjunctions can be folded into Boolean flags. Quantify the output tolerance after this finite proof. In particular we do not induct on ``the limit exists'' or ``this circuit is unitary at all precisions.''
\end{enumerate}
Two details are important for uniformity. A backward pass through a finite rational expression computes all needed precision requests once forward bounds on intermediate magnitudes have been collected. When a variable-length product uses a fixed supplied strong operator $A$, evaluate it by iterating on \emph{finite rational vectors}: approximate $Aq$ to a specified local error from $q$, save the finite result as the next accumulator, and repeat. For example a power of degree $n$ with $\|A\|\le b$ takes per-step error at most $\delta_{k+3}/((n+1)(1+b)^n)$ for absolute output error below $\delta_{k+2}$ given an equally controlled input. This is ordinary numerical primitive recursion. Iteration of a procedure that builds and later evaluates arbitrary function syntax would require a different argument. There is no PR evaluator for all PR programs, by diagonalization. Our schemata compile fixed input programs by substitution. A variable list of gates uses a single joint sampler indexed by the gate number or a fixed finite gate alphabet.

\subsection{What serves as a physical postulate}
Operationally one specifies preparations, admissible effects and an evolution. Their interpretation uses the complex Hilbert-space rules: composition of independent systems is tensor product, isolated reversible time evolution is unitary, and probabilities on compatible outcomes add. These are modeling assumptions about experiments. In finite dimension, assuming a single normalized additive probability assignment on \emph{all effects} already implies the density-matrix trace rule (Theorem~\ref{thm:q-effect}). The choice of Hamiltonian and its parameters is further physical input. We use $\hbar=1$ throughout the dynamics; in physical units the generator is $-iH/\hbar$, with a positive lower bound for $\hbar$ if it is itself named.

An experiment here produces a probability, so its definition does not require a source of random bits. To connect that number with a frequency, take $N$ independent tensor preparations and repeat a two-outcome effect of probability $p$. For the indicator variables $X_j$, tensor factorization gives $\mathbb E X_j=p$ and $\mathbb E X_iX_j=p^2$ for $i\ne j$. Expanding the finite square gives $\mathbb E(\overline X-p)^2=p(1-p)/N\le1/(4N)$. For a rational target $r$ with a certified $|p-r|\le\epsilon/2$, the event $|\overline X-r|\ge\epsilon$ is decidable on its finite outcome strings. On it $(\overline X-p)^2\ge\epsilon^2/4$, hence its total weight is at most $1/(N\epsilon^2)$. An integer $N\ge1/(\alpha\epsilon^2)$ suffices for failure probability $\alpha>0$. This finite-sample assertion also applies when $p$ is an infinite-dimensional Born probability supplied by the later constructions.
